% sec6_7.tex — Section 6.7: Estimation under Conditional Homoskedasticity (Optional)

We saw in Section~3.8 that the GMM estimator reduced to the 2SLS estimator
under conditional homoskedasticity. This section shows how the 2SLS can be
generalized to incorporate serial correlation. It will then become clear that the
GLS estimator is not a GMM estimator with serial correlation. The latter half of
this section examines the relationship between the GLS and GMM.

\subsection*{Kernel-Based Estimation of S under Conditional Homoskedasticity}

The relationship between serial correlation in $\mathbf{g}_t \equiv \mathbf{x}_t \cdot \varepsilon_t$ and serial correlation in
the error term $\varepsilon_t$ becomes clearer under conditional homoskedasticity. Let
\begin{equation}
  \omega_j \equiv \E(\varepsilon_t \varepsilon_{t-j}) \quad \text{for all } j.
  \label{eq:6.7.1}
\end{equation}
Since $\{\varepsilon_t\}$ is stationary by Assumptions~3.1 and~3.2, $\omega_j$ does not depend on $t$.
If $\E(\varepsilon_t) = 0$ (which is the case if $\mathbf{x}_t$ includes the constant), then $\omega_j$ equals the
$j$-th order autocovariance of $\{\varepsilon_t\}$. Suppose the error term is conditionally
homoskedastic:

\medskip
\noindent\textbf{conditional homoskedasticity:}
\begin{equation}
  \E(\varepsilon_t \varepsilon_{t-j} \mid \mathbf{x}_t, \mathbf{x}_{t-j}) = \omega_j
  \quad (j = 0, \pm 1, \pm 2, \ldots).
  \label{eq:6.7.2}
\end{equation}

Under this additional condition, the usual argument utilizing the Law of Total
Expectations can be applied to $\bm{\Gamma}_j$:
\begin{align}
  \bm{\Gamma}_j &= \E(\mathbf{g}_t \mathbf{g}_{t-j}') \quad (\text{since } \E(\mathbf{g}_t) = \mathbf{0}) \notag \\
  &= \E(\varepsilon_t \varepsilon_{t-j} \mathbf{x}_t \mathbf{x}_{t-j}') \notag \\
  &= \E[\E(\varepsilon_t \varepsilon_{t-j} \mid \mathbf{x}_t, \mathbf{x}_{t-j}) \mathbf{x}_t \mathbf{x}_{t-j}']
  \quad \text{(by the Law of Total Expectations)} \notag \\
  &= \omega_j\,\E(\mathbf{x}_t \mathbf{x}_{t-j}').
  \label{eq:6.7.3}
\end{align}
Thus, unless $\E(\mathbf{x}_t \mathbf{x}_{t-j}')$ is zero, $\{\mathbf{g}_t\}$ is serially correlated if and only if $\{\varepsilon_t\}$ is.

To estimate $\bm{\Gamma}_j$, we exploit the fact that $\bm{\Gamma}_j$ is a product of two second moments,
$\E(\mathbf{x}_t \mathbf{x}_{t-j}')$ and $\E(\varepsilon_t \varepsilon_{t-j})$. A natural estimator of $\E(\mathbf{x}_t \mathbf{x}_{t-j}')$ is $\frac{1}{n} \sum_{t=j+1}^{n} \mathbf{x}_t \mathbf{x}_{t-j}'$. It
is easy to show that $\omega_j$ is consistently estimated by $\frac{1}{n} \sum_{t=j+1}^{n} \hat{\varepsilon}_t \hat{\varepsilon}_{t-j}$ where $\hat{\varepsilon}_t$ is
the residual from some consistent estimator (the proof is almost the same as in the
case with no serial correlation, see Proposition~3.2). Thus, a natural estimate of $\bm{\Gamma}_j$
under conditional homoskedasticity is
\begin{equation}
  \hat{\bm{\Gamma}}_j = \left(\frac{1}{n} \sum_{t=j+1}^{n} \hat{\varepsilon}_t \hat{\varepsilon}_{t-j}\right)
  \left(\frac{1}{n} \sum_{t=j+1}^{n} \mathbf{x}_t \mathbf{x}_{t-j}'\right).
  \label{eq:6.7.4}
\end{equation}

Given these estimated autocovariances, the kernel-based estimation of $\hat{\mathbf{S}}$ is exactly
the same as in the case without conditional homoskedasticity described in the previous
section.

\subsection*{Data Matrix Representation of Estimated Long-Run Variance}

For the purpose of relating the GMM estimator with serial correlation but with
conditional homoskedasticity to conventional estimators, the data matrix representation
of $\hat{\mathbf{S}}$ is useful. By defining an $n \times n$ matrix $\hat{\bm{\Omega}}$ suitably, we can write $\hat{\mathbf{S}}$ as
\begin{equation}
  \hat{\mathbf{S}} = \frac{1}{n}\,\mathbf{X}' \hat{\bm{\Omega}} \mathbf{X},
  \label{eq:6.7.5}
\end{equation}
where $\mathbf{X}$ is the $n \times K$ data matrix whose $t$-th row is $\mathbf{x}_t'$. The matrix $\hat{\bm{\Omega}}$ looks much
like the autocovariance matrix of $\{\varepsilon_t\}$:
\begin{equation}
  \underset{(n \times n)}{\hat{\bm{\Omega}}} =
  \begin{bmatrix}
    \hat{\omega}_0 & \hat{\omega}_1 & \hat{\omega}_2 & \cdots & \hat{\omega}_{n-1} \\
    \hat{\omega}_1 & \hat{\omega}_0 & \hat{\omega}_1 & \cdots & \hat{\omega}_{n-2} \\
    \vdots & \ddots & \ddots & \ddots & \vdots \\
    \hat{\omega}_{n-2} & \cdots & \hat{\omega}_1 & \hat{\omega}_0 & \hat{\omega}_1 \\
    \hat{\omega}_{n-1} & \cdots & \hat{\omega}_2 & \hat{\omega}_1 & \hat{\omega}_0
  \end{bmatrix}.
  \label{eq:6.7.6}
\end{equation}
If we know \emph{a priori} that $\Cov(\varepsilon_t, \varepsilon_{t-j}) = 0$ (so that $\bm{\Gamma}_j = \mathbf{0}$) for $j > q$, specify the
elements of this $\hat{\bm{\Omega}}$ as
\begin{equation}
  \hat{\omega}_j = \begin{cases}
    \dfrac{1}{n} \displaystyle\sum_{t=j+1}^{n} \hat{\varepsilon}_t \hat{\varepsilon}_{t-j} & \text{if } 0 \leq j \leq q, \\[8pt]
    0 & \text{if } j > q.
  \end{cases}
  \label{eq:6.7.7}
\end{equation}
Then the $\hat{\mathbf{S}}$ in \eqref{eq:6.6.4} can be written as \eqref{eq:6.7.5}. If we do not know that there is such
a $q$, then use the kernel-based estimator of $\mathbf{S}$. For the Bartlett kernel, let
\begin{equation}
  \hat{\omega}_j = \begin{cases}
    \displaystyle\left[1 - \frac{j}{q(n)}\right] \frac{1}{n} \sum_{t=j+1}^{n} \hat{\varepsilon}_t \hat{\varepsilon}_{t-j}
    & \text{if } 0 \leq j \leq q(n) - 1, \\[8pt]
    0 & \text{if } j > q(n) - 1.
  \end{cases}
  \label{eq:6.7.8}
\end{equation}
Then the consistent estimator \eqref{eq:6.6.3} with \eqref{eq:6.6.4} equals \eqref{eq:6.7.5}.

It is now easy to show the following extension of the 2SLS estimator to encompass
serial correlation:
\begin{itemize}
\item The efficient GMM estimator under conditional homoskedasticity can be written
as
\begin{equation}
  \hat{\bm{\delta}}(\hat{\mathbf{S}}^{-1})
  = [\mathbf{Z}'\mathbf{X}(\mathbf{X}'\hat{\bm{\Omega}}\mathbf{X})^{-1}\mathbf{X}'\mathbf{Z}]^{-1}
  \mathbf{Z}'\mathbf{X}(\mathbf{X}'\hat{\bm{\Omega}}\mathbf{X})^{-1}\mathbf{X}'\mathbf{y},
  \label{eq:6.7.9}
\end{equation}
where $\mathbf{Z}$ and $\mathbf{y}$ are the data matrices of the regressors and the dependent variable.
This generalizes $(3.8.3')$ on page~230.

\item The consistent estimator of $\avar(\hat{\bm{\delta}}(\hat{\mathbf{S}}^{-1}))$ reduces to
\begin{align}
  \widehat{\avar(\hat{\bm{\delta}}(\hat{\mathbf{S}}^{-1}))}
  &= \left[\frac{1}{n}\,\mathbf{Z}'\mathbf{X}
  \Bigl(\frac{1}{n}\,\mathbf{X}'\hat{\bm{\Omega}}\mathbf{X}\Bigr)^{-1}
  \frac{1}{n}\,\mathbf{X}'\mathbf{Z}\right]^{-1} \notag \\
  &= n \cdot [\mathbf{Z}'\mathbf{X}(\mathbf{X}'\hat{\bm{\Omega}}\mathbf{X})^{-1}\mathbf{X}'\mathbf{Z}]^{-1}.
  \label{eq:6.7.10}
\end{align}
So the square roots of the diagonal elements of $[\mathbf{Z}'\mathbf{X}(\mathbf{X}'\hat{\bm{\Omega}}\mathbf{X})^{-1}\mathbf{X}'\mathbf{Z}]^{-1}$ are the standard
errors. This generalizes $(3.8.5')$ on page~230.
\end{itemize}

\subsection*{Relation to GLS}

The procedure for incorporating serial correlation into the GMM estimation we
have described is different from the GLS procedure. To see the difference most
clearly, consider the case where $\mathbf{x}_t = \mathbf{z}_t$. If there are $L$ regressors, $\mathbf{X}'\mathbf{Z}$ and $\mathbf{X}'\hat{\bm{\Omega}}\mathbf{X}$
in \eqref{eq:6.7.9} are all $L \times L$ matrices, so the efficient GMM estimator that exploits the
orthogonality conditions $\E(\mathbf{z}_t \cdot \varepsilon_t) = \mathbf{0}$ is the OLS estimator $\hat{\bm{\delta}}_{\text{OLS}} = (\mathbf{Z}'\mathbf{Z})^{-1}\mathbf{Z}'\mathbf{y}$.
The consistent estimate of its asymptotic variance is obtained by setting $\mathbf{X} = \mathbf{Z}$ in
\eqref{eq:6.7.10}:
\begin{equation}
  \widehat{\avar(\hat{\bm{\delta}}_{\text{OLS}})}
  = n \cdot (\mathbf{Z}'\mathbf{Z})^{-1}(\mathbf{Z}'\hat{\bm{\Omega}}\mathbf{Z})(\mathbf{Z}'\mathbf{Z})^{-1}.
  \label{eq:6.7.11}
\end{equation}

This may seem strange because we know from Section~1.6 that the GLS estimator
is more efficient than OLS when the error is serially correlated. This is a finite
sample result, but, given that we have a consistent estimator $\hat{\bm{\Omega}}$ of $\bm{\Omega}$, we could take
the GLS formula and estimate $\bm{\delta}$ as
\[
  \hat{\bm{\delta}}_{\text{GLS}} = (\mathbf{Z}'\hat{\bm{\Omega}}^{-1}\mathbf{Z})^{-1}\mathbf{Z}'\hat{\bm{\Omega}}^{-1}\mathbf{y}.
\]
Is not this estimator asymptotically more efficient than OLS in that its asymptotic
variance is smaller? The problem with GLS is that consistency is \emph{not} guaranteed,
if the regressors are merely predetermined, as shown below. As we emphasized
in Chapter~2, the regressors are not strictly exogenous in most time series models.
It follows that \emph{GLS should not be used to correct for serial correlation in the
error term for models lacking strict exogeneity}. In particular, for the $\mathbf{x}_t = \mathbf{z}_t$ case,
the correct procedure to adjust for serial correlation is to leave the point estimate
unchanged while incorporating serial correlation in the estimate of the asymptotic
variance.

That the GLS estimator is generally inconsistent can be seen as follows. To
focus on the problem at hand, assume that the true autocovariance matrix is proportional
to some known $n \times n$ matrix $\mathbf{V}$:
\[
  \Var(\varepsilon_1, \varepsilon_2, \ldots, \varepsilon_n) = \sigma^2 \cdot \mathbf{V}.
\]
As shown in Section~1.6, the GLS estimator of the coefficient $\bm{\delta}$ in the regression
model $\mathbf{y} = \mathbf{Z}\bm{\delta} + \bm{\varepsilon}$ can be calculated as the OLS estimator on the transformed model
$\mathbf{C}\mathbf{y} = \mathbf{C}\mathbf{Z}\bm{\delta} + \mathbf{C}\bm{\varepsilon}$, where $\mathbf{C}$ is the ``square root'' of $\mathbf{V}^{-1}$ such that $\mathbf{C}'\mathbf{C} = \mathbf{V}^{-1}$. But
look at the $t$-th observation of the transformed model, which can be written as
\begin{equation}
  \tilde{y}_t = \tilde{\mathbf{z}}_t' \bm{\delta} + \tilde{\varepsilon}_t,
  \label{eq:6.7.12}
\end{equation}
where
\begin{align*}
  \tilde{y}_t &= c_{t1} y_1 + \cdots + c_{tn} y_n, \\
  \tilde{\mathbf{z}}_t &= c_{t1} \mathbf{z}_1 + \cdots + c_{tn} \mathbf{z}_n, \\
  \tilde{\varepsilon}_t &= c_{t1} \varepsilon_1 + \cdots + c_{tn} \varepsilon_n,
\end{align*}
$(c_{t1}, \ldots, c_{tn})$ is the $t$-th row of $\mathbf{C}$.
For the GLS estimator to be consistent, it is necessary that $\E(\tilde{\mathbf{z}}_t \cdot \tilde{\varepsilon}_t) = \mathbf{0}$. The
left-hand side of this condition can be written as
\[
  \E(\tilde{\mathbf{z}}_t \cdot \tilde{\varepsilon}_t) = \text{(a)} + \text{(b)},
  \quad \text{(a)} = \sum_{s=1}^{n} c_{ts}^2\,\E(\mathbf{z}_s \cdot \varepsilon_s),
  \quad \text{(b)} = \sum_{s \neq v} c_{ts} \cdot c_{tv}\,\E(\mathbf{z}_s \cdot \varepsilon_v).
\]
By the orthogonality condition $\E(\mathbf{z}_t \cdot \varepsilon_t) = \mathbf{0}$, (a) is zero. For~(b), because
regressors are merely predetermined and not assumed to be strictly exogenous,
$\E(\mathbf{z}_s \cdot \varepsilon_v)$ is not guaranteed to be zero for $s \neq v$. So~(b) is not zero, unless the $c$'s
take particular values to offset the nonzero terms.

There is one important special case where GLS is consistent, and that is when
the error is a finite-order autoregressive process. To take the simplest case, suppose
that $\{\varepsilon_t\}$ is AR(1): $\varepsilon_t = \phi \varepsilon_{t-1} + \eta_t$. Its autocovariance matrix is
\[
  \Var(\varepsilon_1, \varepsilon_2, \ldots, \varepsilon_n) = \sigma^2 \cdot \mathbf{V}, \quad
  \mathbf{V} = \frac{1}{1 - \phi^2}
  \begin{bmatrix}
    1 & \phi & \phi^2 & \cdots & \phi^{n-1} \\
    \phi & 1 & \phi & \cdots & \phi^{n-2} \\
    & & \ddots & & \\
    \phi^{n-2} & \cdots & \phi & 1 & \phi \\
    \phi^{n-1} & \cdots & \phi^2 & \phi & 1
  \end{bmatrix}.
\]
It is easy to verify that the square root, $\mathbf{C}$, of $\mathbf{V}^{-1}$ (satisfying $\mathbf{C}'\mathbf{C} = \mathbf{V}$) is given by
\begin{equation}
  \mathbf{C} =
  \begin{bmatrix}
    \sqrt{1 - \phi^2} & 0 & 0 & \cdots & 0 & 0 \\
    -\phi & 1 & 0 & \cdots & 0 & 0 \\
    0 & -\phi & 1 & \cdots & 0 & 0 \\
    \vdots & \vdots & \vdots & & \vdots & \vdots \\
    0 & 0 & 0 & \cdots & -\phi & 1
  \end{bmatrix}.
  \label{eq:6.7.13}
\end{equation}
So the $t$-th transformed equation is
\begin{equation}
  y_t - \phi y_{t-1} = (\mathbf{z}_t - \phi \mathbf{z}_{t-1})' \bm{\delta} + \eta_t.
  \label{eq:6.7.14}
\end{equation}
The GLS estimate of $\bm{\delta}$, which is the OLS estimate in the regression of $y_t - \phi y_{t-1}$
on $\mathbf{z}_t - \phi \mathbf{z}_{t-1}$, is consistent if $\eta_t$ is orthogonal to \emph{both} $\mathbf{z}_t$ and $\mathbf{z}_{t-1}$. However, those
orthogonality conditions are different from the orthogonality conditions for the
untransformed model, which are that $\E(\mathbf{z}_t \cdot \varepsilon_t) = \mathbf{0}$.

\subsection*{Questions for Review}

\begin{enumerate}
\item Verify \eqref{eq:6.7.5} for $n = 3$, $q = 1$.

\item (Sargan's statistic) \; Derive the generalization of Sargan's statistic $(3.8.10')$ on
page~230 to the case of serially correlated errors.

\item Suppose that $\mathbf{z}_t$ is a strict subset of $\mathbf{x}_t$. We have seen in Section~3.8 that the
2SLS estimator reduces to the OLS estimator. Is this true here as well? That
is, does \eqref{eq:6.7.9} reduce to $(\mathbf{Z}'\mathbf{Z})^{-1}\mathbf{Z}'\mathbf{y}$? [Answer: No.]

\item (Data matrix representation of $\hat{\mathbf{S}}$ without conditional homoskedasticity) \; Define
$\hat{\bm{\Omega}}$ suitably so that \eqref{eq:6.7.5} equals the Bartlett kernel-based estimator of $\hat{\mathbf{S}}$
without conditional homoskedasticity, namely, \eqref{eq:6.6.5} with \eqref{eq:6.6.7}. \textbf{Hint:} Write
$\hat{\bm{\Omega}}$ as \eqref{eq:6.7.6}. If the lag length $q$ is known, $\hat{\omega}_j = \hat{\varepsilon}_j \hat{\varepsilon}_{t-j}$ for $0 \leq j \leq q$ and 0 for
$j > q$. How would you define $\hat{\omega}_j$ for the Bartlett kernel-based estimator?
\end{enumerate}
